\section*{الفصل \ref{ch:b2:weather-systems} --- الأنظمة الجوية}

\begin{solution}{exo:b2:weather-systems:1}
(أ) $f = 2 \times 7.292 \times 10^{-5} \times \sin\ang{50} = \qty{1.117e-4}{s^{-1}}$. (ب) مع
$\bar T = \qty{278}{K}$، يكون $\tan\alpha = 1.117 \times 10^{-4} \times 278 \times 15/(9.81 \times 6)
= 0.00791$، أي نحو 1 من 126. (ج) $1/0.00791 = 126$: على بعد نحو \qty{126}{km} خلف الجبهة. (د) الميل
يتناسب عكسيًّا مع فرق الحرارة: فينخفض إلى النصف، إلى 1 من 253.
\end{solution}

\begin{solution}{exo:b2:weather-systems:2}
(أ) $6 \times 150 = \qty{900}{km}$ أمامها. (ب) $900/25 = \qty{36}{h}$، أي يوم ونصف. (ج) أوّلًا
السمحاق العالي، ثم السمحاق الطبقي، ثم الطبقي المتوسط؛ فالسحب تتكاثف وتنخفض، ويبدأ المطر
من السحب المنخفضة الكثيفة كلّما اقتربت الجبهة.
\end{solution}

\begin{solution}{exo:b2:weather-systems:3}
تشير الرموز إلى اتجاه حركة الجبهة، وإلى الهواء الأدفأ أمام الجبهة الباردة أو
الهواء الأبرد أمام الجبهة الدافئة. (أ) جبهة باردة: يقع الهواء البارد إلى شمالها الغربي، خلفها؛
وتتحرّك نحو الجنوب الشرقي. (ب) جبهة دافئة: يقع الهواء البارد إلى شمالها، أمامها؛ وتتحرّك
شمالًا. (ج) جبهة مسدودة، تصل المنخفض بالنقطة الثلاثية التي تبدأ منها الجبهتان الباردة والدافئة؛
ويقع الهواء البارد على جانبيها، والهواء الدافئ مرفوع فوقه؛ وتتحرّك في الاتجاه الذي تشير إليه
رموزها. (د) جبهة ثابتة: تشير المثلّثات جنوبًا نحو الهواء الدافئ، فيقع الهواء البارد
إلى الشمال؛ وتكاد لا تتحرّك.
\end{solution}

\begin{solution}{exo:b2:weather-systems:4}
(أ) $\beta = 2 \times 7.292 \times 10^{-5} \times \cos\ang{45}/(6.371 \times 10^6) =
\qty{1.62e-11}{m^{-1}.s^{-1}}$. (ب) للموجة التي طولها \qty{4000}{km}، $k = \qty{1.571e-6}{m^{-1}}$، و$\beta/k^2 =
\qty{6.6}{m/s}$ و$c = 12 - 6.6 = \qty{5.4}{m/s}$؛ وللتي طولها \qty{8000}{km}، $\beta/k^2 =
\qty{26.3}{m/s}$ و$c = \qty{-14.3}{m/s}$. (ج) تتحرّك الموجة القصيرة شرقًا، \qty{470}{km} في اليوم؛
وتتحرّك الطويلة غربًا، \qty{1240}{km} في اليوم، عكس الريح.
\end{solution}

\begin{solution}{exo:b2:weather-systems:5}
(أ) $c = 0$ حين يكون $U = \beta/k^2$، فيكون $k = \sqrt{\beta/U}$ و$\lambda = 2\pi/k =
2\pi\sqrt{U/\beta}$. (ب) $2\pi\sqrt{12/(1.62 \times 10^{-11})} = \qty{5410}{km}$. (ج) يبلغ طول دائرة
العرض \ang{45} القيمة $2\pi \times 6371 \times \cos\ang{45} = \qty{28300}{km}$: أي نحو 5.2
موجة، فخمس موجات عمليًّا. (د) يزداد $\lambda$ بالشكل $\sqrt U$: فالريح الغربية الأقوى تجعل
الموجات الثابتة أطول، فيتّسع نصف الكرة لعدد أقلّ منها.
\end{solution}

\begin{solution}{exo:b2:weather-systems:6}
(أ) عند $C_+ = (\sqrt{72}, \sqrt{72}, 27)$: $\dot x = 10(\sqrt{72} - \sqrt{72}) = 0$، $\dot y =
\sqrt{72}(28 - 27) - \sqrt{72} = 0$، $\dot z = 72 - \tfrac83 \times 27 = 72 - 72 = 0$.
(ب) $\mathrm{e}^{-13.67} = 1.16 \times 10^{-6}$: يتقلّص الحجم نحو مليون مرّة في وحدة
زمنية واحدة. (ج) يمكن أن يتقلّص حجم بينما يتمدّد في اتجاه واحد وينضغط بشدّة أكبر بكثير
في الاتجاهات الأخرى: فتنمو المسافة على امتداد الاتجاه المتمدّد، ويتباعد مساران متجاوران
على امتداده، بينما يتلاشى الحجم بينهما.
\end{solution}

\begin{solution}{exo:b2:weather-systems:7}
(أ) $\ln 2/0.906 = 0.765$ وحدة زمنية. (ب) $\ln(10/10^{-8})/0.906 = \ln(10^9)/0.906 = 22.9$ وحدة
زمنية. (ج) $\ln 1000/0.906 = 7.6$ وحدة زمنية: فتحسين الحالة الابتدائية ألف مرّة
لا يشتري إلا ثلث الزمن الذي وُجد في (ب)، لأن $\ln 10^3/\ln 10^9 = 1/3$.
\end{solution}

\begin{solution}{exo:b2:weather-systems:8}
(أ) $(303 - 195)/195 = 0.554$ و$V_{\max} = \sqrt{0.9 \times 0.554 \times 15\,000} =
\qty{86.5}{m/s}$. (ب) أعلى بكثير من \qty{70}{m/s} للفئة 5، وقريبة من \qty{86}{m/s} التي
تجاوزتها خمس عواصف حديثة. (ج) نادرًا ما تبقى العواصف طويلًا في الحالة المستقرّة التي تفترضها
النظرية؛ والقصّ الرأسي للريح، والهواء الجاف المسحوب إلى العاصفة والهواء الهابط من حولها، وهي ما
تُغفله النظرية، تُضعفها.
\end{solution}

\begin{solution}{exo:b2:weather-systems:9}
(أ) $(303 - 195)/303 = 0.356$. (ب) $(303 - 195)/195 = 0.554$؛ والنسبة $0.554/0.356 =
303/195 = 1.554$. (ج) في محرّك كارنو يغادر الشغلُ المحرّك. أمّا في الإعصار المداري فيتبدّد
الشغل بالاحتكاك قرب سطح البحر، وتلك الحرارة تدفّئ الهواء المتدفّق نحو الداخل عند
درجة الحرارة $T_s$، حيث يأخذ المحرّك حرارته: فتُستعمل ثانية، ويسلك المحرّك
كأن مردوده أكبر بالعامل $T_s/T_o$.
\end{solution}

\begin{solution}{exo:b2:weather-systems:10}
(أ) اشتقاق \cref{prop:b2:weather-systems:slope} دون تقريب يعطي
$\mathrm{d}z/\mathrm{d}x = f(\rho_c v_c - \rho_w v_w)/[g(\rho_c - \rho_w)]$. والمائع الأكثف
المستقرّ فوق مائع أخفّ غير مستقرّ: فأدنى اضطراب يجعله ينقلب. ففي جبهة
ساكنة يجب إذن أن يقع الهواء البارد الأكثف في الأسفل، وهو الإسفين الوارد في التعريف؛ وعندئذ
$\rho_c > \rho_w$، ويرتفع السطح نحو الهواء البارد، و$\rho_c v_c > \rho_w v_w$.
(ب) $\zeta = \partial v/\partial x = 10/(50 \times 10^3) = \qty{2e-4}{s^{-1}}$، أي نحو $1.9f$ مع
$f = \qty{1.03e-4}{s^{-1}}$. (ج) موجبة في نصف الكرة الشمالي: إعصارية، وبقدر
معامل كوريوليس نفسه؛ فالجبهات أشرطة من الدوّامية الإعصارية الشديدة.
\end{solution}

\begin{solution}{exo:b2:weather-systems:11}
(أ) $\log_2(6/0.5) = \log_2 12 = 3.58$ تضاعف، و$3.58 \times 2.4 = 8.6$~يوم. (ب) كلّ
يوم إضافي من التنبّؤ يتطلّب نموًّا أقلّ بزمن تضاعف واحد قبل أن يصبح الخطأ عديم الجدوى: فيجب أن
يتقلّص الخطأ الابتدائي بالعامل $2^{1/2.4} = 1.33$ كلّ عقد. (ج) التخفيض بعشر مرّات يساوي
$\log_2 10 = 3.32$ تضاعف؛ فإن لم يكسب إلا 4 إلى 5 أيام، فالأخطاء الصغيرة تتضاعف في $4/3.32 = 1.2$
إلى $5/3.32 = 1.5$ يوم: فالأخطاء الأصغر تنمو أسرع من الأكبر التي قيست في
تنبّؤات ثمانينيات القرن العشرين، عبر الحرارة الكامنة المتحرّرة في الحمل على المقاييس الصغيرة.
\end{solution}

\begin{solution}{exo:b2:weather-systems:12}
(أ) عند \qty{30}{\celsius}، $e_s = \qty{43.5}{hPa}$ و$q_s^* = \qty{27.1}{g/kg}$؛ وللهواء عند
\qty{29}{\celsius} القيمة $e = 0.8 \times 41.0 = \qty{32.8}{hPa}$ و$q_a = \qty{20.4}{g/kg}$؛
و$h_o^* - h^* = 1004.5 \times 1 + 2.501 \times 10^6 \times 0.0067 = \qty{17.8}{kJ/kg}$، مقابل
\qty{16.8}{kJ/kg}. (ب) $(303.15 - 200)/200 = 0.516$ و$V_{\max} = \sqrt{0.9 \times 0.516
\times 17\,800} = \qty{90.9}{m/s}$، مقابل \qty{87.9}{m/s}. (ج) نحو \qty{3.4}{\percent} لكل
كلفن. فعدم التوازن يرتفع بنسبة \qty{6}{\percent}، فيضيف \qty{3}{\percent} إلى السرعة؛
ويرتفع العامل $(T_s - T_o)/T_o$ بنسبة \qty{1}{\percent}، فيضيف \qty{0.5}{\percent}: فالوقود،
أي الرطوبة التي يستطيع البحر الأدفأ تقديمها، هو الأكثر إسهامًا، عبر الارتفاع الأسّي
للضغط $e_s$.
\end{solution}

\begin{solution}{pb:b2:weather-systems:1}
\textbf{1.} $(480 - 390)/6 = \qty{15}{km/h}$.
\textbf{2.} $390/15 = \qty{26}{h}$ بعد الساعة 12 بالتوقيت العالمي المنسّق: أي الساعة 14 من اليوم التالي.
\textbf{3.} $f = \qty{1.117e-4}{s^{-1}}$.
\textbf{4.} مع $\bar T = \qty{281}{K}$، يكون $\tan\alpha = 1.117 \times 10^{-4} \times 281 \times
16/(9.81 \times 8) = 0.0064$، أي نحو 1 من 156.
\textbf{5.} تكون الجبهة عندئذ على بعد $6 \times 15 = \qty{90}{km}$ وراء البلدة؛ ويقع السطح على ارتفاع
$90\,000 \times 0.0064 = \qty{576}{m}$ فوقها.
\textbf{6.} $1500/0.0064 = \qty{234}{km}$ خلف الجبهة، ويُبلغ بعد $234/15 = \qty{15.6}{h}$ من
المرور.
\textbf{7.} سحب ركامية وركامية مزنية، وزخّات، وربّما رعد، وهبّات ريح وهبوط مفاجئ في
الحرارة؛ والضغط، الذي ينخفض مع اقتراب الجبهة، يبدأ بالارتفاع حين تمرّ؛ ثم
هواء أبرد وأصفى.
\textbf{8.} $\beta = 2 \times 7.292 \times 10^{-5} \times \cos\ang{50}/(6.371 \times 10^6) =
\qty{1.47e-11}{m^{-1}.s^{-1}}$.
\textbf{9.} $k = 2\pi/(5 \times 10^6) = \qty{1.257e-6}{m^{-1}}$ و$\beta/k^2 = \qty{9.3}{m/s}$.
\textbf{10.} $c = 20 - 9.3 = \qty{10.7}{m/s}$، نحو الشرق.
\textbf{11.} $10.7 \times 86\,400 = \qty{924}{km}$ نحو الشرق.
\textbf{12.} $\lambda_s = 2\pi\sqrt{20/(1.47 \times 10^{-11})} = \qty{7330}{km}$؛ ويبلغ طول دائرة
العرض \ang{50} القيمة $2\pi \times 6371 \times \cos\ang{50} = \qty{25700}{km}$، أي نحو 3.5
موجة كهذه.
\textbf{13.} أمامه، على جانبه الشرقي، حيث تُحدث الأمواج القصيرة هواءً صاعدًا
وينخفض الضغط السطحي.
\textbf{14.} $\lambda = \ln(3.9 \times 10^{-3}/5 \times 10^{-7})/(14 - 4) = \ln 7800/10 = 0.896$.
\textbf{15.} $\ln 2/0.896 = 0.77$ وحدة زمنية.
\textbf{16.} أقلّ بنسبة \qty{1}{\percent} من 0.906: متوسط 200 زوج، وزمن محدود، وقراءات من
رسم بياني.
\textbf{17.} متى تباعد التشغيلان كثيرًا، وقعا في نقطتين لا علاقة بينهما من الجاذب،
ولا يمكن أن تتجاوز المسافة بينهما حجمه: فيتشبّع النموّ.
\textbf{18.} بعد $\ln 10/0.896 = 2.6$ وحدة زمنية.
\textbf{19.} $\log_2(3/0.1) = \log_2 30 = 4.91$ تضاعف.
\textbf{20.} $4.91 \times 2.4 = 11.8$~يوم.
\textbf{21.} زمن تضاعف واحد: 2.4~يوم.
\textbf{22.} $4.91 \times 5 = 24.5$~يوم: كان التقدير الأوّل متفائلًا أكثر بكثير ممّا ينبغي.
\textbf{23.} مع زمن تضاعف قدره 2.4 يوم، يكسب خطأ أصغر بعشر مرّات $3.32 \times 2.4
= 8.0$~يوم، لا 4 إلى 5: فالأخطاء الصغيرة تنمو أسرع من تضاعف كلّ 2.4 يوم، والتقدير
المبنيّ على زمن التضاعف ذاك يبالغ في الأفق.
\textbf{24.} يمكن أن يكون هذا التنبّؤ مفيدًا \textbf{12 يومًا} على الأكثر تقريبًا، وفي الواقع
أقلّ من ذلك.
\end{solution}
